The direct-prefix formula leaves a reflected cutoff and one terminal prefix.
Both admit positive-kernel recurrences, and together they close the replacement
derivative ledger.

\begin{lemma}[Direct positive-prefix areas beyond the endpoint; proved]
\label{lem:terminal-modal-static-direct-lobes-D8}
For every $D\ge8$, the direct branch obeys
\begin{equation}
 \sum_{K\ge1}(\Pi_8^\circ(K))_+=\frac{311}{1024},
 \qquad
 \sum_{K\ge1}(\Pi_D^\circ(K))_+
 \le\frac{356}{1024}<\frac{29}{64}\quad(D\ge9),
 \qquad
 \sum_{K\ge1}(\Pi_D^\circ(K))_-<\frac{57}{8}.
 \label{eq:terminal-modal-static-direct-lobes-D8}
\end{equation}
\end{lemma}

\begin{proof}
For $D\ge6$, $K\ge3$, coefficient extraction in
\eqref{eq:terminal-modal-static-direct-prefix-gf} simplifies to
\begin{equation}
 \Pi_D^\circ(K)=2^{-K-2}[z^D](z^2+z-6)(1+z)^K
 (z+2+z^{K+2}).
 \label{eq:terminal-modal-static-direct-prefix-simple}
\end{equation}
Split the last factor into the main term $z+2$ and alias term $z^{K+2}$,
writing the corresponding coefficients as $M_D(K)$ and $A_D(K)$.  If
$r=D-2$ and $x=K-D+3$, then
\begin{equation}
 \sum_{x\ge0}M_D(D-3+x)y^x
 =\frac{4-7y^2-3y^3}{(2-y)^{r+3}}.
 \label{eq:terminal-modal-static-direct-main-gf}
\end{equation}
Increasing $D$ by one convolves this coefficient sequence with the
nonnegative geometric kernel $(2-y)^{-1}$, whose coefficient mass is one.
Therefore its positive-part mass cannot increase.  Exact extraction at the
first two distances gives
\begin{equation}
 \sum_K(M_8(K))_+=\frac{279}{1024},
 \qquad
 \sum_K(M_9(K))_+=\frac{57}{256}.
 \label{eq:terminal-modal-static-direct-main-base}
\end{equation}

For the alias put $u=D-K-4$.  Its exact coefficient is
\begin{equation}
 A_D(K)=2^{-K-2}\left\{
 \binom Ku+\binom K{u+1}-6\binom K{u+2}\right\}.
 \label{eq:terminal-modal-static-direct-alias}
\end{equation}
Pascal's identity gives the exact smoothing relation
\begin{equation}
 A_{D+2}(K+1)=\frac12\{A_D(K)+A_{D+1}(K)\}.
 \label{eq:terminal-modal-static-direct-alias-smoothing}
\end{equation}
The positive-part masses at $D=8,9$ are respectively $1/16$ and $1/8$.
Induction in \eqref{eq:terminal-modal-static-direct-alias-smoothing} bounds
every later one by $1/8$.  Direct extraction at $D=8$ gives total positive
area $311/1024$.  For $D\ge9$, $(x+y)_+\le x_++y_+$ and the main sequence is
bounded by its $D=9$ value, giving
$57/256+1/8=356/1024$.  This proves the first bound in
\eqref{eq:terminal-modal-static-direct-lobes-D8}.  Finally
\eqref{eq:terminal-modal-static-Q-moments} and summation by parts give
negative mass equal to positive mass plus
$20/3+(4/3)(-1/2)^D$.  For $D\ge8$ the latter is at most $427/64$, so
\[
 \frac{356}{1024}+\frac{427}{64}
 =\frac{7188}{1024}<\frac{57}{8},
\]
proving the second direct bound.
\end{proof}

\begin{lemma}[Reflected positive area; proved]
\label{lem:terminal-modal-static-reflected-prefix}
For $m\ge64$, $N=m-2$, $8\le D\le m$, and $H=2m-D$,
\begin{equation}
 \sum_{K=1}^{N-1}(\widehat\Pi_H(K))_+<\frac1{1024},
 \qquad
 \widehat\Pi_H(K)=\sum_{k=1}^K\widehat Q_{m,D}(k).
 \label{eq:terminal-modal-static-reflected-positive}
\end{equation}
\end{lemma}

\begin{proof}
For $H-k\ge2$, direct use of
\eqref{eq:terminal-modal-static-a-atom} gives the coefficientwise increment
identity
\begin{equation}
 \widehat Q_H(k)=[z^H]2^{-k}z^{k-1}(1+z)^{k-1}
 \left(\frac14-z^2\right)(z^2+2z-3).
 \label{eq:terminal-modal-static-reflected-increment-gf}
\end{equation}
For $H>4$ and $K\le H-2$, summing these coefficients in $k$ gives, with
\[
 B(z)=\frac{z(1+z)}2,
 \qquad
 A(z)=\frac{(z+3)(z^2-1/4)}{z+2},
\]
the coefficientwise cone identity
\begin{equation}
 \widehat\Pi_H(K)=[z^H]A(z)\{1-B(z)^K\}.
 \label{eq:terminal-modal-static-reflected-prefix-closed}
\end{equation}
The polynomial $A(1-B)$ has degree four.  Hence, whenever $H>4$ and
$K\le H-4$,
\begin{equation}
 \widehat\Pi_H(K+1)
 =\frac12\{\widehat\Pi_{H-1}(K)+\widehat\Pi_{H-2}(K)\}.
 \label{eq:terminal-modal-static-reflected-prefix-recurrence}
\end{equation}
Put
\[
 \mathfrak A_H=\sum_{K=1}^{H-3}(\widehat\Pi_H(K))_+,
 \qquad
 e_H=(\widehat\Pi_H(H-2))_+.
\]
Convexity of the positive part and
\eqref{eq:terminal-modal-static-reflected-prefix-recurrence} give
\begin{equation}
 \mathfrak A_H\le
 \frac12\mathfrak A_{H-1}+\frac12\mathfrak A_{H-2}
 +\frac12e_{H-2}.
 \label{eq:terminal-modal-static-reflected-area-recurrence}
\end{equation}
Exact dyadic extraction gives
\begin{equation}
 \mathfrak A_{64}=\frac{8330122021020985}{2^{63}},
 \qquad
 \mathfrak A_{65}=\frac{14501951681351529}{2^{64}},
 \label{eq:terminal-modal-static-reflected-area-bases}
\end{equation}
and the sole boundary error has the closed form
\begin{equation}
 e_H=\frac{6H^2-32H-7+15(-1)^H}{2^{H+3}}>0
 \qquad(H\ge64).
 \label{eq:terminal-modal-static-reflected-boundary}
\end{equation}
The elementary identity
$\sum_{H\ge M}H^2/2^H=2^{1-M}(M^2+2M+3)$ gives
\[
 \sum_{H\ge64}e_H<\frac{12681}{2^{65}}
 <\frac1{1024}-\mathfrak A_{64}
 =\frac{677077233720007}{2^{63}}.
\]
Induction in \eqref{eq:terminal-modal-static-reflected-area-recurrence}
therefore proves $\mathfrak A_H<1/1024$ for every $H\ge64$.
The literal range $K\le m-3$ is a subrange of $K\le H-3$ because
$H=2m-D\ge m$.  This proves the claim, including the central case $H=m$.
\end{proof}

\begin{lemma}[Folded terminal prefix; proved]
\label{lem:terminal-modal-static-terminal-prefix}
For the same range,
\begin{equation}
 |\Pi_{m,D}(N)|<\frac12.
 \label{eq:terminal-modal-static-terminal-prefix}
\end{equation}
\end{lemma}

\begin{proof}
For $D\le N+1$, reverse the reflected coefficient in
\eqref{eq:terminal-modal-static-reflected-prefix-closed} and combine it with
\eqref{eq:terminal-modal-static-direct-prefix-simple}.  Exact extraction is
\begin{equation}
 \Pi_{m,D}(N)=2^{-m}[x^D](1+x)^N\mathcal V(x),
 \label{eq:terminal-modal-static-terminal-prefix-gf}
\end{equation}
where
\begin{align}
 \mathcal V(x)
 &=(x-2)(x+2)(x+3)
 -\frac{x^2(1+3x)(4-x^2)}{1+2x}\notag\\
 &=-\frac{399}{32}-\frac{49}{16}x-\frac{23}{8}x^2
 +\frac34x^3+\frac32x^4+\frac{15}{32(1+2x)}.
 \label{eq:terminal-modal-static-terminal-prefix-partial}
\end{align}
Let $p_i=2^{-N}\binom Ni$, $p_\star=\max_i p_i$, and
\[
 A_D^N=[x^D]\frac{(1+x)^N}{1+2x}.
\]
Pascal induction gives $0\le A_D^N\le\binom ND$ for $D\le N-1$, while
$A_N^N=(-1)^N$ and $A_{N+1}^N=-2(-1)^N$.  Consequently
\[
 -\frac{151}{32}p_\star\le\Pi_{m,D}(N)
 \le\frac{87}{128}p_\star.
\]
The central binomial bound
\[
 p_\star\le2^{-62}\binom{62}{31}\le\frac1{\sqrt{94}}
\]
follows from
$4^{-n}\binom{2n}{n}\le(3n+1)^{-1/2}$ and monotonicity between adjacent
parities.  Since $151^2<256\cdot94$, both sides have absolute value less
than $1/2$.

At the remaining central boundary $D=N+2=m$, the direct alias contains its
exceptional coincident atom.  Literal extraction gives
\[
 \Pi_{m,m}(N)=2^{-m}
 \left\{\frac34N^2-\frac{71}{8}+\frac{15}{8}(-1)^N\right\},
\]
which lies in $(0,1/2)$ for $N\ge62$.
\end{proof}

\begin{lemma}[Folded weighted tail; proved]
\label{lem:terminal-modal-static-folded-tail}
For $m\ge64$, $N=m-2$, and $8\le D\le m$, put
\[
 \mathcal E_{m,D}=\sum_{K=1}^{N-1}\Pi_{m,D}(K)
 +\frac{20}{3}+\frac43(-1/2)^D.
\]
Then $\mathcal E_{m,D}>0$.
\end{lemma}

\begin{proof}
For $D\le N+1$, pairing the omitted direct tail with the reflected prefix in
\eqref{eq:terminal-modal-static-direct-prefix-gf} and
\eqref{eq:terminal-modal-static-reflected-prefix-closed} gives
\begin{equation}
 \mathcal E_{m,D}=2^{-(m-1)}[x^D](1+x)^N\mathcal R(x),
 \label{eq:terminal-modal-static-folded-tail-gf}
\end{equation}
where
\begin{align}
 \mathcal R(x)
 &=\frac{12+52x+61x^2+3x^3-20x^4-16x^5+x^6+3x^7}
 {(1-x)(1+2x)^2}\notag\\
 &=\frac{113}{64}+\frac{37}{8}x+\frac{55}{16}x^2
 -\frac14x^3-\frac34x^4+\frac{32}{3(1-x)}
 -\frac{49}{96(1+2x)}+\frac{5}{64(1+2x)^2}.
 \label{eq:terminal-modal-static-folded-tail-partial}
\end{align}
Besides $A_D^N$ above, put
\[
 B_D^N=[x^D]\frac{(1+x)^N}{(1+2x)^2}.
\]
Multiplication by $1+x$ gives Pascal recurrences in $N$.
The boundary $A_{N-1}^N=(1-(-1)^N)/2$ proves
\[
 0\le A_D^N\le\binom ND\qquad(0\le D\le N-1).
\]
The corresponding boundary for $B$ is zero when $N$ is even and
$(N-1)/2$ when $N$ is odd, proving
$B_D^N\ge0$ for $0\le D\le N-3$.

For $D\le N-3$, retain the positive prefix term $32/[3(1-x)]$ in
\eqref{eq:terminal-modal-static-folded-tail-partial}, use these two bounds,
and discard all other positive terms.  The resulting lower bound is
\[
 \frac{32}{3}\sum_{r=0}^D\binom Nr
 -\frac14\binom N{D-3}-\frac34\binom N{D-4}
 -\frac{49}{96}\binom ND>0,
\]
because each of the three displayed negative terms is charged to its own
occurrence in the prefix sum.
For the last four indices $D=N-2,N-1,N,N+1$, the exact boundary values are
\begin{align*}
 A_D^N&=\left(\lfloor N/2\rfloor,
 \frac{1-(-1)^N}{2},(-1)^N,-2(-1)^N\right),\\
 B_D^N&=\left((-1)^N\lfloor N/2\rfloor,
 (-1)^{N+1}N,(-1)^N(2N+1),(-1)^{N+1}4(N+1)\right).
\end{align*}
Let $S_D=\sum_{r=0}^D\binom Nr$.  Charge the two negative binomial terms to
their occurrences in $S_D$, leaving $(29/3)S_D$.  The displayed boundary
values give
\[
 \frac{29}{3}S_D-\frac{49}{96}|A_D^N|-\frac5{64}|B_D^N|
 >\frac{29}{3}(2^N-N-1)-N>0.
\]
This proves
\eqref{eq:terminal-modal-static-folded-tail-gf} is positive on $D\le N+1$.

At the central distance $D=N+2=m$, the same literal atom correction as in
\cref{lem:terminal-modal-static-terminal-prefix} gives instead
\begin{equation}
 \mathcal E_{m,m}=2^{-(m-1)}[x^{N+2}](1+x)^N\mathcal R(x)
 +\frac6{2^m}.
 \label{eq:terminal-modal-static-folded-tail-central}
\end{equation}
Here
\[
 A_{N+2}^N=4(-1)^N,
 \qquad B_{N+2}^N=(-1)^N(8N+12).
\]
The same charge leaves a lower bound
\[
 \frac{29}{3}2^N-\frac{49}{24}-\frac5{64}(8N+12)>0
\]
for the coefficient in
\eqref{eq:terminal-modal-static-folded-tail-central}.  The additional central
correction is positive as well, completing the proof.
\end{proof}

\begin{theorem}[Replacement derivative bound; proved]
\label{thm:terminal-modal-static-derivative-replacement}
For every $m\ge64$ and $0\le j\le m-6$, the derivative term in
\eqref{eq:terminal-modal-static-folded-source-ledger} obeys
\[
 \mathcal V_{m,m-j}\ge-\frac{1421}{14400}q>-\frac q8.
\]
Consequently the interior two-step cone
\eqref{eq:terminal-modal-static-two-step-cone} holds strictly on these rows.
Only its five direct frontier rows remain open.
\end{theorem}

\begin{proof}
The $D=6,7$ prefix bounds are proved in
\cref{lem:terminal-modal-static-derivative-prefix-gf}.  For $D\ge8$,
\eqref{eq:terminal-modal-static-direct-lobes-D8} and
\eqref{eq:terminal-modal-static-reflected-positive} give
\[
 \sum_{K=1}^{N-1}(\Pi_{m,D}(K))_+
 \le\frac{357}{1024}<\frac{29}{64}<\frac12.
\]
If $P_+$ and $P_-$ are the two folded prefix areas, then
\cref{lem:terminal-modal-static-folded-tail} gives
\[
 P_-<P_++\frac{427}{64}
 \le\frac{357}{1024}+\frac{427}{64}
 =\frac{7189}{1024}<\frac{57}{8}.
\]
The endpoint estimate is
\eqref{eq:terminal-modal-static-terminal-prefix}.  Thus all conditions
\eqref{eq:terminal-modal-static-replacement-lobes} hold, and
\cref{lem:terminal-modal-static-derivative-replacement} proves the derivative
bound.  Combining it with the mass and base bounds gives the strict margin
\eqref{eq:terminal-modal-static-replacement-total-margin}, proving the
interior cone.
\end{proof}
